% Problem Setup - ICML style

\paragraph{Road graph and way-level routes.}
The environment is a directed road graph $G=(V,E)$ extracted from OpenStreetMap, where each node $v \in V$ has geographic coordinates and each edge $e \in E$ carries a road type (motorway, primary, secondary, residential, service) and length.
An OSM \emph{way} is a maximal sequence of edges forming a single named, typed road segment---the semantic atom of road network topology.
A ground-truth route is a way sequence $\mathbf{w} = (w_1, \dots, w_M)$ obtained by map-matching a GPS trajectory to $G$ and performing consecutive deduplication.
Way-level tokenization compresses route length from $L \sim$ hundreds of nodes to $M \sim$ tens of ways, while preserving topological adjacency: consecutive ways share at least one node.

\paragraph{Conditioning.}
Each route is conditioned on $c = (o, d, t, x)$: origin~$o$, destination~$d$, departure time~$t$ (from which cyclical features such as hour-of-day and day-of-week are derived), and optional spatial context~$x$.

\paragraph{Desiderata.}
A route generator $p_\theta(\mathbf{w} \mid c)$ should satisfy three properties:
\begin{enumerate}
  \item \textbf{Intent consistency}: the generated route reaches the destination~$d$.
  \item \textbf{Graph feasibility}: consecutive ways share a node, so the output is a valid path on $G$.
  \item \textbf{Corridor multi-modality}: for OD pairs with structurally distinct alternatives, the generator should cover the empirical distribution of corridors.
\end{enumerate}

\paragraph{Evaluation metrics.}
Let $\mathbf{w}$ be a ground-truth route and $\{\hat{\mathbf{w}}^{(k)}\}_{k=1}^K$ be $K$ generated routes for the same conditioning~$c$.
We evaluate four complementary aspects:

\textit{Arrival Rate} is the fraction of generated routes that successfully reach~$d$ within a maximum step budget.

\textit{MeanMaxJaccard} measures how well the best generated route matches the ground truth.
Let $W(\mathbf{w})$ denote the set of ways traversed by route~$\mathbf{w}$.
For each ground-truth route, we compute the best Jaccard match among $K$ candidates:
\begin{equation}
  J_{\max} = \max_{k \in \{1,\dots,K\}} \frac{|W(\hat{\mathbf{w}}^{(k)}) \cap W(\mathbf{w})|}{|W(\hat{\mathbf{w}}^{(k)}) \cup W(\mathbf{w})|}.
\end{equation}
MeanMaxJaccard is the average of $J_{\max}$ across all test routes.

\textit{Coverage-$\tau$ AUC} captures how many distinct ground-truth corridor patterns are covered by the generated set, without committing to a single similarity threshold.
For a given OD group with multiple ground-truth routes, Coverage@$K(\tau)$ is the fraction of ground-truth routes for which at least one of the $K$ generated routes achieves $J \geq \tau$.
The threshold-free aggregate is:
\begin{equation}
  \mathrm{CovAUC} = \frac{1}{\tau_{\max} - \tau_{\min}} \int_{\tau_{\min}}^{\tau_{\max}} \mathrm{Coverage}@K(\tau) \, d\tau,
  \label{eq:covauc}
\end{equation}
averaged across OD groups.
This metric penalizes methods that cover only a single corridor mode regardless of how well they match it.

\textit{Self-Diversity@$K$} is the average pairwise dissimilarity (1 $-$ Jaccard) among the $K$ generated routes for the same conditioning, measuring how distinct the generated alternatives are.
